% Physical and virtual cuts.

\centering
\begin{tikzpicture}[setA, x=1cm, y=1cm]

% Physical cut.
\begin{scope}[xshift=0cm]

\node[font=\sffamily\footnotesize\bfseries, anchor=south]
  at (3.25, 3.95) {Physical cut};
\node[font=\sffamily\scriptsize\itshape, text=gMid, anchor=south]
  at (3.25, 3.50) {what DBLog \emph{does not} produce};

\draw[axisLine] (0.10, 0) -- (6.40, 0);

\def\xF{3.25}

% Counterfactual snapshot with the same target state as the virtual cut.
% Centered at y=2.30, the same height as the right panel's equality unit.
\node[blockFill, anchor=center, align=center,
      font=\sffamily\scriptsize, inner xsep=5pt, inner ysep=3pt]
  (lstate) at (\xF, 2.30) {%
  all keys read at $f$\\[1pt]
  \begin{tabular}{@{\ }r@{\,:\,}l@{\ }}
    100 & $3000$  \\
    200 & $12000$ \\
    300 & $3500$  \\
  \end{tabular}};

% Single cut line at f. The 0.9pt weight matches the dependency arrows in Figure 1.
\draw[gDark, line width=0.9pt] (\xF, 0) -- (lstate.south);
\node[font=\sffamily\scriptsize\bfseries, anchor=north, inner sep=1pt]
  at (\xF, -0.05) {$f$};

% Subtitle at y=-0.90, aligned with the right panel's subtitle.
\node[font=\sffamily\scriptsize\itshape, text=gDark, anchor=north,
      align=center, text width=58mm]
  at (\xF, -0.90)
  {a single source-log coordinate at which\\
   every key in scope is observed};

% Annotation at y=0.55, aligned with the right panel's annotation.
\node[font=\sffamily\scriptsize, anchor=west, text=gMid,
      align=left, inner sep=1pt]
  at (\xF+0.20, 0.55)
  {monolithic\\\strut single-line cut};

\end{scope}

\draw[gMid, line width=0.5pt] (7.30, -1.30) -- (7.30, 4.10);

% Virtual cut.
\begin{scope}[xshift=7.95cm]

\node[font=\sffamily\footnotesize\bfseries, anchor=south]
  at (3.25, 3.95) {Virtual cut};
\node[font=\sffamily\scriptsize\itshape, text=gMid, anchor=south]
  at (3.25, 3.50) {what DBLog \emph{does} produce};

\draw[axisLine] (0.10, 0) -- (6.40, 0);
\def\xCa{0.95}
\def\xCb{2.45}
\def\xCc{3.95}
\def\xCd{5.50}
\foreach \x/\lab in {\xCa/c_1, \xCb/c_2, \xCc/c_3} {
  \draw[axisTick] (\x, 0.06) -- (\x, -0.06);
  \node[font=\sffamily\scriptsize, anchor=north, inner sep=1pt]
    at (\x, -0.06) {$\lab$};
}
\draw[axisTick, line width=0.8pt] (\xCd, 0.10) -- (\xCd, -0.10);
\node[font=\sffamily\scriptsize\bfseries, anchor=north, inner sep=1pt]
  at (\xCd, -0.08) {$c_4 = f$};

\draw[braceMirror] (\xCa-0.30, -0.35) -- (\xCa+0.30, -0.35);
\node[font=\sffamily\scriptsize, anchor=north, inner sep=1pt]
  at (\xCa, -0.45) {chunk $A$};
\draw[braceMirror] (\xCc-0.30, -0.35) -- (\xCc+0.30, -0.35);
\node[font=\sffamily\scriptsize, anchor=north, inner sep=1pt]
  at (\xCc, -0.45) {chunk $B$};

\draw[guide] (\xCd, 0) -- (\xCd, 1.95);

% Replay/source equality on the claimed keys.
\def\eqY{2.30}
\node[font=\sffamily\Large, anchor=center] (eq) at (3.25, \eqY) {$=$};
\node[blockFill, anchor=east, align=center,
      font=\sffamily\scriptsize, left=3pt of eq.west,
      inner xsep=5pt, inner ysep=3pt]
  (lbox) {%
  \textbf{$\Apply(\CleanPrefix) \restrict K$}\\[1pt]
  \begin{tabular}{@{\ }r@{\,:\,}l@{\ }}
    100 & $3000$  \\
    200 & $12000$ \\
    300 & $3500$  \\
  \end{tabular}};
\node[block, anchor=west, align=center,
      font=\sffamily\scriptsize, right=3pt of eq.east,
      inner xsep=5pt, inner ysep=3pt]
  (rbox) {%
  \textbf{$\Src(b_0, H, f) \restrict K$}\\[1pt]
  \begin{tabular}{@{\ }r@{\,:\,}l@{\ }}
    100 & $3000$  \\
    200 & $12000$ \\
    300 & $3500$  \\
  \end{tabular}};

\node[font=\sffamily\scriptsize, anchor=center, text=gMid,
      align=center, inner sep=1pt]
  at (3.25, 0.55)
  {\emph{no single line}\\ outcomes match at $f$};

\node[font=\sffamily\scriptsize\itshape, text=gDark, anchor=north,
      align=center, text width=58mm]
  at (3.25, -0.90)
  {chunk reads at distinct coordinates\\
   replay-at-frontier $=$ source-at-frontier on scope $K$};

\end{scope}

\end{tikzpicture}

\caption{Physical cut vs.\ virtual cut. \emph{Left:} a physical cut
requires a single source-log coordinate at which every key in scope is
observed, shown as a single vertical line. DBLog does \emph{not} produce
one. The values shown are the counterfactual outcome on $K$. \emph{Right:}
a virtual cut requires replay at the frontier to equal the source at
the frontier on the chosen scope. The chunk reads (as in
Figure~\ref{fig:running-example-timeline}) sit at distinct earlier
coordinates, but the replay reconstructs the source state at the frontier.}
\label{fig:virtual-vs-physical-cut}
